% ============================================================================
% PATCH v5.7.01b: Add Curvature Fork Section
% ============================================================================
% Add a dedicated section about the curvature fork and category errors
% ============================================================================

\section{The Curvature Fork: Avoiding Category Errors}

A critical issue in previous formulations of information-based physics has been \textbf{category errors}—conflating mathematical objects from different domains. The L-ToEC framework explicitly distinguishes three types of ``curvature'' that appear at different layers:

\subsection{Three Distinct Curvature Objects}

\begin{enumerate}
    \item \textbf{Physical Spacetime Curvature ($\Cphys$)}:
    \[
    \Cphys = \{R^a_{bcd}, R_{ab}, R, \text{and their invariants}\}
    \]
    These are geometric objects on the spacetime manifold $\mathcal{M}$ in layer $\interface$ (L3). They describe gravitational effects.
    
    \item \textbf{Informational/Statistical Curvature ($\Cinfo$)}:
    \[
    \Cinfo = \{R_F(\theta), \Gamma^F_{ij,k}(\theta), \text{Fisher-Rao metric invariants}\}
    \]
    These live on statistical manifolds $(\Theta, g_F)$ in layers $\protocol$ (L1) and $\logic$ (L2), where $g_F$ is the Fisher metric.
    
    \item \textbf{Topos Truth-Value Curvature ($\Ctopos$)}:
    \[
    \Ctopos = \{\text{invariants of } \Omega \text{ in observer topos } \mathcal{E}\}
    \]
    These are logical/computational invariants in the subobject classifier of an observer's topos.
\end{enumerate}

\subsection{The Fundamental Distinction}

\begin{tcolorbox}[colback=red!5!white,colframe=red!75!black,title=Critical Notation Rule]
\[
\Cphys \neq \Cinfo \neq \Ctopos
\]
These objects inhabit different mathematical categories:
\begin{itemize}
    \item $\Cphys$: Riemannian geometry (smooth manifolds)
    \item $\Cinfo$: Information geometry (statistical manifolds) 
    \item $\Ctopos$: Topos theory (Heyting algebras)
\end{itemize}
\textbf{Any map between them must be explicitly defined as a functor with stated properties.}
\end{tcolorbox}

\subsection{Explicit Functor Definitions (When They Exist)}

\paragraph{Physical→Informational Map}
A map $F_{\mathrm{phys\to info}}: \Cphys \to \Cinfo$ exists only when:
\begin{enumerate}
    \item A statistical model $\rho(x|\theta)$ is specified, where $x$ are spacetime measurements.
    \item The Fisher metric $g_F(\theta)$ is computed from $\rho(x|\theta)$.
    \item The map's properties (functoriality, continuity, etc.) are verified.
\end{enumerate}

\paragraph{Informational→Topos Map}  
A map $F_{\mathrm{info\to topos}}: \Cinfo \to \Ctopos$ exists only when:
\begin{enumerate}
    \item An observer topos $\mathcal{E}$ is defined (site, topology).
    \item A proposition language $\mathcal{P}$ is specified.
    \item A valuation $V: \mathcal{P} \to \Gamma(\Omega)$ is given.
    \item The map from statistical states to propositions is defined.
\end{enumerate}

\subsection{Consequences for Claims About ``Curvature → Qualia''}

Any claim of the form ``curvature gives rise to qualia'' must specify:
\begin{enumerate}
    \item Which curvature: $\Cphys$, $\Cinfo$, or $\Ctopos$?
    \item Which explicit functor: $F_{\mathrm{phys\to info}}$, $F_{\mathrm{info\to topos}}$, or a composition?
    \item What are the functor's domain, codomain, and properties?
\end{enumerate}

\textbf{Unspecified mappings are category errors.} Previous formulations that used physical curvature $R_{\mathrm{phys}}$ in discussions of consciousness without specifying the functor $F_{\mathrm{phys\to topos}}$ were committing this error.

\subsection{Implementation in the Framework}

In L-ToEC v5.7+, we enforce:
\begin{itemize}
    \item \textbf{Explicit notation:} Always use $\Cphys$, $\Cinfo$, $\Ctopos$ to distinguish.
    \item \textbf{Functor declarations:} Any map between categories must be declared.
    \item \textbf{Audit trail:} The dependency DAG tracks which functors are used.
    \item \textbf{Z3 verification:} Automated checking that category boundaries are respected.
\end{itemize}

This formalization eliminates the single biggest technical liability identified in the critique.
